\section{割线法、收敛阶与混合策略}
\label{sec:08-Numerical-Analysis/02-Root-Finding/03}

一个黑盒函数，每次调用要跑几分钟仿真，或者花掉一笔真金白银。你要求它的根。上一节的 Newton 法每步需要导数，可你没有：函数不提供导数，自己实现一遍又和重写整个模型差不多。二分法不要导数，但每步只买 \(0.3\) 个十进制位。

中间地带存在吗？不要导数，又比二分快得多。

\subsection{用两点连线冒充切线}

Newton 步是 \(x_{k+1}=x_k-f(x_k)/f'(x_k)\)。没有 \(f'(x_k)\) 时最省事的补救，是用最近两次函数值算一个差商顶上：

\[
f'(x_k)\ \approx\ \frac{f(x_k)-f(x_{k-1})}{x_k-x_{k-1}} .
\]

代入 Newton 步，得到割线法（secant method）

\[
x_{k+1}=x_k-f(x_k)\,\frac{x_k-x_{k-1}}{f(x_k)-f(x_{k-1})} .
\]

几何上，它把过 \((x_{k-1},f(x_{k-1}))\) 与 \((x_k,f(x_k))\) 的那条割线当成曲线的替身，取它与横轴的交点。和 Newton 法是同一个念头，只是切线换成了割线。

\begin{center}
\begin{tikzpicture}[x=1.95cm,y=0.32cm,font=\small,>=stealth]
  \draw[->,black!55] (1.05,0) -- (3.75,0) node[right,font=\footnotesize] {\(x\)};
  \draw[blue!75,thick] plot[domain=1.15:3.55,samples=70] (\x,{\x*\x-2});
  \draw[red!70,thick] (1.8067,0) -- (3.55,10.46);
  \draw[black!55,dashed] (1.6846,0) -- (3.2,7.88);
  \fill[black!75] (3.4,9.56) circle (1.4pt);
  \fill[black!75] (2.6,4.76) circle (1.4pt);
  \node[left,font=\footnotesize] at (3.32,10.1) {\(x_{k-1}\)};
  \node[left,font=\footnotesize] at (2.52,5.4) {\(x_k\)};
  \draw[black!45,dashed] (1.4142,0.25) -- (1.4142,4.3);
  \node[above,font=\footnotesize,black!60] at (1.4142,4.3) {根};
  \fill[red!70] (1.8067,0) circle (1.4pt);
  \draw[black!55] (1.6846,0) circle (1.4pt);
  \draw[black!40] (1.98,-0.9) -- (1.83,-0.25);
  \node[right,font=\footnotesize,red!70] at (1.95,-1.5) {割线给出的 \(x_{k+1}\)};
  \draw[black!40] (1.30,-0.9) -- (1.66,-0.25);
  \node[left,font=\footnotesize,black!60] at (1.33,-1.5) {切线本会给出的点};
\end{tikzpicture}

\emph{曲线仍是 \(f(x)=x^2-2\)。红色实线是过两个已知点的割线，虚线是 \(x_k\) 处的真切线。割线的落点 \(1.807\) 比切线的落点 \(1.685\) 离根远一些，因为差商只是导数的近似；但它一次导数也没用，全部信息来自两个函数值。}
\end{center}

成本结构是这个方法真正漂亮的地方：每步只需要一次新的函数计算。上一步算过的 \(f(x_{k-1})\) 缓存下来，这一步直接复用，于是每个函数值服务了两步迭代——当前步的分子和下一步的分母。用有限差分近似导数再跑 Newton 就没有这个便宜：那样每步要花两次函数计算才凑出一个斜率。

代价也藏在同一个公式里。分母是 \(f(x_k)-f(x_{k-1})\)，两点靠得很近时它是两个相近数相减，相对误差被放大，这正是浮点相消，见\hyperref[sec:08-Numerical-Analysis/01-Floating-Point-and-Error/02]{``误差如何产生与放大''}。分母一小，步长就大，一步跳到哪里没人能保证。

\subsection{黄金比例：比线性快，比二次慢}

在足够光滑函数的简单根附近，割线法的误差满足

\[
e_{k+1}\approx C\,e_k\,e_{k-1},
\]

其中 \(e_k=x_k-r\)。来历是两点线性插值的余项公式——割线就是过两点的一次插值多项式，余项里恰好带着 \((x-x_{k-1})(x-x_k)\) 这个因子，两个误差的乘积由此而来。这里只用它的结论。

递推式的形状已经把答案写好了一半：当前误差是上两步误差之积，比线性收敛的 \(e_{k+1}\approx\lambda e_k\) 强，又比 Newton 的 \(e_{k+1}\approx Ce_k^2\) 弱——毕竟 \(e_{k-1}\) 总比 \(e_k\) 大。把它变成一个数，只需要一个渐近假设：设收敛后期 \(\abs{e_k}\sim\abs{e_{k-1}}^p\)。代进去，

\[
\abs{e_{k-1}}^{p^2}\sim\abs{e_{k+1}}\approx\abs{C}\,\abs{e_k}\,\abs{e_{k-1}}\sim\abs{e_{k-1}}^{\,p+1},
\]

两边指数相等给出 \(p^2=p+1\)，正根是黄金比例

\[
p=\varphi=\frac{1+\sqrt5}{2}\approx1.618 .
\]

超线性，但够不着二次。

按每步计算，Newton 更快。按每次函数调用计算，账要重算一遍：Newton 的一步同时消耗一次函数与一次导数，若两者成本相当，折合每次调用的效率因子是 \(\sqrt2\approx1.414\)；割线法一步只花一次函数调用，效率因子就是 \(1.618\)。函数昂贵而导数不便宜时，割线法反而占优。反过来，若导数由自动微分几乎免费获得——这在今天的数值代码里越来越常见——天平立刻倒向 Newton。收敛阶不能脱离信息成本单独比较。

\subsection{收敛阶只描述最后一段}

先把比较的语言统一。若 \(x_k\to r\)，且存在 \(p\ge1\) 与有限的 \(\lambda>0\) 使

\[
\lim_{k\to\infty}\frac{\abs{e_{k+1}}}{\abs{e_k}^{\,p}}=\lambda,
\]

就称该序列具有 Q-阶 \(p\)。\(p=1\) 时还要求 \(\lambda<1\) 才算线性收敛，\(p>1\) 称超线性，\(p=2\) 称二次。按这个尺子，二分法是 \(p=1,\ \lambda=1/2\)；收缩不动点迭代是 \(p=1\)，\(\lambda\) 就是不动点处导数的绝对值；简单根附近的割线法约 \(1.618\)，Newton 法是 \(2\)。

麻烦在于定义里那个 \(k\to\infty\)。它描述的只是最后一段。

\begin{center}
\begin{tikzpicture}[x=0.66cm,y=0.20cm,font=\small,>=stealth]
  \draw[->,black!55] (0,0) -- (13.2,0) node[right,font=\footnotesize] {迭代步 \(k\)};
  \draw[->,black!55] (0,0.6) -- (0,-18.8);
  \foreach \y in {-4,-8,-12,-16}
    \draw[black!15] (0,\y) -- (12.6,\y);
  \node[left,font=\footnotesize] at (-0.2,-4) {\(-4\)};
  \node[left,font=\footnotesize] at (-0.2,-8) {\(-8\)};
  \node[left,font=\footnotesize] at (-0.2,-12) {\(-12\)};
  \node[left,font=\footnotesize] at (-0.2,-16) {\(-16\)};
  \draw[black!45,dashed] (0,-16) -- (12.6,-16);
  \node[above,font=\footnotesize,black!55] at (2.6,-15.8) {舍入平台};
  \draw[black!30,dashed] (3.6,0.4) -- (3.6,-17.6);
  \draw[blue!75,thick] (0,-0.3) -- (12,-3.91);
  \draw[black!70,thick] (0,-0.3) -- (1,-0.42) -- (2,-0.58) -- (3,-0.85) -- (4,-1.3)
    -- (5,-2.05) -- (6,-3.3) -- (7,-5.3) -- (8,-8.5) -- (9,-13.6) -- (10,-16.2)
    -- (11,-16.0) -- (12,-16.3);
  \draw[red!75,thick] (0,-0.3) -- (1,-0.45) -- (2,-0.65) -- (3,-1.0) -- (4,-1.8)
    -- (5,-3.5) -- (6,-7.0) -- (7,-14.0) -- (8,-16.1) -- (9,-15.9) -- (10,-16.2)
    -- (11,-16.0) -- (12,-16.3);
  \node[right,font=\footnotesize,blue!70] at (12.1,-3.9) {二分法};
  \node[left,font=\footnotesize,red!75] at (5.9,-8.4) {Newton};
  \node[right,font=\footnotesize,black!70] at (9.3,-11.4) {割线法};
  \node[font=\footnotesize,black!55] at (1.8,-18.3) {全局段};
  \node[font=\footnotesize,black!55] at (7.6,-18.3) {渐近段与平台段};
  \node[left,font=\footnotesize] at (-0.2,0.6) {\(\log_{10}\abs{e_k}\)};
\end{tikzpicture}

\emph{三条曲线在前三四步几乎重合：这一段由函数的全局形状主导，谁也快不起来。分隔线之后才轮到收敛阶说话，Newton 与割线的曲线陡然下坠，二分法仍是那条不变的斜线。最后所有曲线都撞在同一条水平线上——舍入平台，误差不再下降。收敛阶只解释中间那一段。}
\end{center}

于是``实测收敛阶''这件事就变得可疑。真根已知时可以观察 \(\abs{e_{k+1}}/\abs{e_k}^p\) 是否趋于常数；真根未知时通常拿相邻迭代差顶替误差，用三四个连续点反解 \(p\)。但在图上的第一段，拟合出来的数字描述的是全局风景而非根的局部结构；在最后一段，\(x_k\) 与 \(x_{k+1}\) 只是两个相邻浮点数，比值可能是 \(1\)、\(0.5\)、\(2\) 这类毫无几何含义的数。此外，振荡型方法、多步方法、有限精度下提前终止的序列，未必具有简单 Q-阶，硬套定义会得到看起来很精确的误导。

经验阶靠得住的用法只有一种：在你已经有理由相信自己进了渐近区之后，用它做一次确认。反过来用——靠它判断收敛是否已经发生——不成立。

\subsection{割线点不受任何约束，所以要配括区}

比收敛阶更直白的问题是安全性。割线法不保证下一步落在含根区间里。即便起手的两个端点异号，后续迭代点照样可能飞到几百个单位之外：函数含噪声、附近有近乎水平的平台、或者还藏着另一个根，斜率估计都会被扭曲，而扭曲的斜率直接变成扭曲的步长。

工程上的成熟做法是混合。始终维护一个异号括区 \([a,b]\)，这是从二分法继承的安全网；每步先尝试割线步，或者用三点做逆二次插值得到一个更好的候选；然后审查这个候选——落在括区之外、进展小到像在原地踏步、或者分母危险地接近零，一律丢弃，改用一次二分步；最后用新点更新括区。

Brent 方法（及其变体 Brent--Dekker）是这套思路的标杆实现，也是许多科学计算库里 \texttt{brentq} 之类函数的内核。它的逻辑很朴素：绝大多数步由插值完成，快；偶尔插一步二分，保证最坏情况下区间仍在稳定收缩。二分法提供的下界与插值提供的上界，在这里被同时拿到。

\subsection{选方法时到底在比较什么}

收敛阶只回答一个问题：一切顺利时它有多快。选方法的智力活儿在别处。

首先问能不能可靠地找到异号括区。能找到，混合方法就有了安全底牌；找不到，就得选自带全局化策略的方案，比如阻尼 Newton 或信赖域。其次问函数本身：连续吗，光滑吗，带噪声吗——含噪函数上的割线斜率可能被一个离群的采样值整个劫持，而二分法只看符号，对幅值噪声不敏感。再问导数的成本：自动微分改变了这道题的答案，导数几乎免费时 Newton 与拟 Newton 的优势会大幅扩张。最后问需求本身：要一张误差证书，还是要一个快速近似；是单次求根，还是同一个模型要解上百万次——后者值得花成本做预处理或构造更好的初值。

一句可以带走的判断：求根失败的代价越高，越应该把安全网放在速度前面。嵌在另一个优化循环内部的求根尤其如此，外层不会告诉你内层悄悄发散了，你只会看到一堆莫名其妙的结果。

三节走下来，方法的谱系已经清楚：二分法只用符号，慢而不失手；Newton 法用切线，快而无保障；割线法用两点连线，居中而便宜。它们共享同一个模式——先用一个容易解的模型替换难解的函数，再解模型。二分法的模型是``符号''，Newton 的模型是切线，割线法的模型是过两点的一次插值。下一单元把这个替换动作单独拿出来研究：给定若干数据点，用多项式或样条去逼近未知函数，误差从哪里来、什么时候会失控，见\hyperref[ch:063]{``插值与逼近：通过数据点，不等于理解点之间''}。
